\documentclass[12pt]{article}
\usepackage{amsmath, amssymb}
\usepackage[a4paper, margin=1in]{geometry}

\begin{document}

\title{Lecture 23: Predicting the Impossible – Tail Bounds and System Reliability}
\date{}
\maketitle

(Coding Session – Tails and Bounds)

\section*{1. Introduction (Slide 1)}
Title: Predicting the Impossible: Tail Bounds and System Reliability

Context: Case Study: Amazon Black Friday Sale 2026

Focus: Probability bounds, uncertainty, and system reliability

\section*{2. Amazon Black Friday Sale 2026 (Slide 2)}

Timeline (as given):
\begin{itemize}
\item Expected start: Friday, November 27, 2026
\item Early deals: November 20–21, 2026
\item Main sale period: November 25–27, 2026
\item Cyber Monday: November 30, 2026
\end{itemize}

Target Matrix (as shown):
\begin{itemize}
\item Electronics \& Tech → Discount ceiling: 75\%
\item Amazon Devices → Discount ceiling: 50\%
\item Fashion \& Apparel → Discount ceiling: 60\%
\item Home \& Kitchen → Discount ceiling: 60\%
\end{itemize}

Additional factors:

Amazon Prime advantages:
\begin{itemize}
\item Early access to inventory
\item Exclusive cashback offers
\end{itemize}

Partner bank leverage:
\begin{itemize}
\item 10\% instant discount
\item 5\% multiplier cashback
\end{itemize}

\section*{3. Connecting with Uncertainty and Computing (Slide 3)}

System Modeling

Multiple servers: Each server processes jobs

Condition of interest: Highlighted server has $\ge k$ jobs

Poisson Arrivals (Datacenter)

Job arrivals modeled as a Poisson process

Time window: 1 hour

Event of interest: Many arrivals within one hour

Condition:

\section*{4. Tail Probability Objective (Slide 4)}

Goal:

Interpretation: Probability that system load exceeds threshold

Context: Used to predict:
\begin{itemize}
\item Server crashes
\item System overload
\end{itemize}

\section*{5. Mean vs Tail Behavior (Slide 5)}

Statement: The mean is safe

The tail is catastrophic

Simulation Details:
\begin{itemize}
\item Server traffic: 1000 minutes
\item Distribution: Poisson with $\lambda = 50$
\end{itemize}

Number of arrivals $\ge k$

\[
P(X \ge k)
\]

Threshold: $k = 60$

Observation:

Crash occurs when: $X \ge 60$

Empirical crash rate:

$87/1000 = 8.7\%$

Required computation:

$P(X \ge 60)$

\section*{6. Hierarchy of Bounds (Slide 6)}

Statement: Better inputs yield tighter bounds

Level 1: Markov’s Inequality (Mean)

Level 2: Chebyshev’s Inequality (Variance)

Level 3: Chernoff Bound (Full distribution, MGF)

Insight:

Exact probability may not be computable

Upper bounds are used

More information $\Rightarrow$ less over-estimation

\section*{7. Markov’s Inequality (Slide 7)}

\[
P(X \ge a) \le \frac{E[X]}{a}
\]

Interpretation: Requires only mean

Decay Behavior: Linear decay

Observation: Over-estimates significantly

\section*{8. Chebyshev’s Inequality (Slide 8)}

\[
P(|X - \mu| \ge a) \le \frac{\sigma^2}{a^2}
\]

Equivalent form:
\[
P(X \ge a) \le \frac{\sigma^2}{(a - \mu)^2}
\]

Decay Behavior: Polynomial decay

Observation: Tighter than Markov

Still over-estimates

\section*{9. Chernoff Bound (Slide 9)}

\[
P(X \ge a) \le \min_{t>0} e^{-ta} M_X(t)
\]

Where:
\[
M_X(t) = E[e^{tX}]
\]

Alternative form:
\[
P(X \ge a) \le e^{-at + \ln M_X(t)}
\]

Decay Behavior: Exponential decay

Observation: Tightest bound

Closely matches exact probability

\section*{10. Comparison of Decay (Slide 10)}

Markov: $1/k$

Chebyshev: $1/k^2$

Chernoff: $e^{-k}$

Conclusion: Chernoff closely maps reality

Markov remains pessimistic

\section*{11. Case Study: SLA Breach (Slide 11)}

Distribution: Poisson

Mean: $\lambda = 500$

Threshold: $k = 600$

Objective:
\[
P(X \ge 600)
\]

Task:
\begin{itemize}
\item Predict probability
\item Size backup capacity
\end{itemize}

\section*{12. Cost of Loose Bound (Slide 12)}

Exact risk: $\sim 1.4\%$

Chernoff: $\sim 8.8\%$

Chebyshev: $\sim 20\%$

Markov: $\sim 90.9\%$

Interpretation:

Markov predicts extremely high failure probability

Leads to unnecessary provisioning

\section*{13. Financial Impact (Slide 13)}

\[
\text{Waste Multiplier} = \frac{\text{Exact Probability}}{\text{Bound Estimate}}
\]

At $\lambda = 500$:

Chernoff: $\sim 10.7\times$

Chebyshev: $\sim 6422\times$

Markov: $\sim 107039\times$

Conclusion:

Loose bounds $\Rightarrow$ massive cost

\section*{14. Over-estimation Interpretation (Slide 14)}

Statement:

Mathematical over-estimation = financial over-provisioning

Explanation:

Over-estimated probability $\Rightarrow$ more backup servers

Leads to unused capacity

\section*{15. Different Distributions (Slide 15)}

Cases:
\begin{itemize}
\item Normal → stable jobs
\item Poisson → independent requests
\item Binomial → fixed users
\item Gamma → bursty traffic
\end{itemize}

Observation:

Chernoff works well for:
\begin{itemize}
\item Poisson
\item Binomial
\end{itemize}

Limitation:

For heavy-tailed (Gamma): Over-estimation occurs

Alternative methods needed

\section*{16. Decision Matrix (Slide 16)}

\begin{tabular}{|c|c|c|c|}
\hline
Method & Requires & Decay & Use Case \\
\hline
Markov & Mean & $1/k$ & Rough check \\
Chebyshev & $\mu, \sigma^2$ & $1/k^2$ & Unknown distribution \\
Chernoff & Full MGF & $e^{-k}$ & Precise sizing \\
\hline
\end{tabular}

\section*{17. Laplace Distribution Task (Slide 17)}

\[
X \sim \text{Laplace}(\mu, b)
\]

Properties:
\begin{itemize}
\item Sharp central peak
\item Heavy tails
\end{itemize}

Threshold:
\[
a = \mu + 3\sigma
\]

Relation:
\[
\sigma = b\sqrt{2}
\]

Objective: Derive Chernoff bound

\section*{18. Derivation Blueprint (Slide 18)}

Step 1: Compute MGF
\[
M_X(t) = E[e^{tX}]
\]

Step 2: Apply Chernoff Inequality
\[
P(X \ge a) \le e^{-ta} M_X(t)
\]

Step 3: Optimize

Differentiate w.r.t $t$

Set derivative = 0

Solve for optimal $t$

Step 4: Substitute

Plug optimal $t$ into bound

\section*{End of Lecture Scribe}

\end{document}